\documentclass[11pt]{article}
\usepackage[margin=1in]{geometry}
\usepackage{amsmath, amssymb}
\usepackage{setspace}
\usepackage{booktabs}

\setstretch{1.2}
\setlength{\parindent}{0pt}

\begin{document}

\begin{center}
{\Large \textbf{Solutions -- Practice Problem Set 1}}\\
\vspace{1em}
{\large ECON 204 -- Contemporary Economic Issues}\\
\vspace{1em}
{\large Winter 2026}
\end{center}
\hrule
\vspace{1em}

%====================================================
\section*{Part I: Inflation, Cost of Living, and Monetary Policy}

\subsection*{Question 1: Measuring Inflation}

CPI last year \(=152.0\), CPI this year \(=158.1\).

\subsubsection*{(1) Inflation rate}

\[
\pi=\frac{CPI_t-CPI_{t-1}}{CPI_{t-1}}\times 100
=\frac{158.1-152.0}{152.0}\times 100
\]

Compute the numerator:
\[
158.1-152.0=6.1
\]

Now divide:
\[
\frac{6.1}{152.0}\approx 0.0401316
\]

Convert to percent:
\[
0.0401316\times 100\approx 4.01\%
\]

\textbf{Answer:} Inflation \(\approx 4.01\%\).

\subsubsection*{(2) Real wage change}

Nominal wage growth \(=3.5\%\). Approximate real wage growth:
\[
\Delta w^{real}\approx \Delta w^{nominal}-\pi
=3.5\%-4.01\%=-0.51\%
\]

\textbf{Answer:} Real wage falls by about \(0.5\%\).

\subsubsection*{(3) Purchasing power interpretation}

If inflation exceeds nominal wage growth, the worker’s pay buys fewer goods and services than before. Even with a raise, the worker is effectively worse off in real terms.
\newpage
\subsection*{Question 2: Cost of Living and Household Budgets}

Budget shares: housing 30\%, food 20\%, transport 15\%, utilities 10\%, other 25\%. Price changes: housing 8\%, food 5\%, transport 4\%, utilities 6\%, other 2\%.

\subsubsection*{(1) Household-specific inflation rate}

This is a weighted average:
\[
\pi^{HH}=\sum w_i\pi_i
\]

Convert shares to decimals:

Housing \(w_H=0.30\), \(\pi_H=0.08\)

Food \(w_F=0.20\), \(\pi_F=0.05\)

Transport \(w_T=0.15\), \(\pi_T=0.04\)

Utilities \(w_U=0.10\), \(\pi_U=0.06\)

Other \(w_O=0.25\), \(\pi_O=0.02\)
\vspace{0.1cm}\\

Compute each contribution:

Housing: \(0.30\times 0.08=0.024\)

Food: \(0.20\times 0.05=0.010\)

Transport: \(0.15\times 0.04=0.006\)

Utilities: \(0.10\times 0.06=0.006\)

Other: \(0.25\times 0.02=0.005\)

Sum:
\[
0.024+0.010+0.006+0.006+0.005=0.051
\]

Convert to percent:
\[
0.051\times 100=5.1\%
\]

\textbf{Answer:} Household inflation \(5.1\%\).

\subsubsection*{(2) Why differs from CPI}

The CPI is based on an “average” basket. A household’s basket differs because:

Different expenditure weights (e.g., renters vs homeowners; families vs singles).

Different local prices (housing and transport are highly regional).

Some households are more exposed to items with high inflation (like rent).

\subsubsection*{(3) Who most affected when housing rises fastest}

Households with high housing shares: renters, low- and middle-income households, single parents, and households in tight metro rental markets. Since housing is a necessity, they have less ability to substitute away.
\newpage
\subsection*{Question 3: Real Interest Rates and Monetary Policy}

Nominal interest rate \(i=5.5\%\). Expected inflation \(\pi^e=2.8\%\).

\subsubsection*{(1) Real interest rate}

Using Fisher approximation:
\[
r\approx i-\pi^e=5.5\%-2.8\%=2.7\%
\]

\textbf{Answer:} \(2.7\%\).

\subsubsection*{(2) How real rate affects decisions}

A higher real interest rate:

Raises the true cost of borrowing \(\rightarrow\) discourages consumption financed by credit (cars, housing) and discourages business investment.

Increases the return to saving \(\rightarrow\) encourages saving.

Generally reduces aggregate demand over time.

\subsubsection*{(3) How a policy rate hike lowers inflation (channels)}

Textbook transmission mechanism:

Interest rate channel: borrowing costs rise \(\rightarrow\) consumption and investment fall.

Housing channel: mortgages more expensive \(\rightarrow\) housing demand slows \(\rightarrow\) house price growth moderates; also reduces related spending.

Exchange rate channel (Canada): higher rates can appreciate CAD \(\rightarrow\) imports cheaper \(\rightarrow\) lowers inflation.

Expectations/credibility: credible tightening anchors inflation expectations \(\rightarrow\) wage/price setting becomes less inflationary.

\subsection*{Question 4: Inflation Expectations}

Expected inflation rises from 2\% to 5\%.

\subsubsection*{(1) Wage bargaining}

Workers seek higher nominal wages to maintain purchasing power. They may request wage increases closer to expected inflation plus productivity growth.

\subsubsection*{(2) Firm response}

Firms anticipate higher wage costs and may raise prices preemptively, especially in sectors with pricing power.

\subsubsection*{(3) Expectations-driven persistence}

If wages rise because inflation is expected to rise, and prices rise because wages rose, expectations can become self-fulfilling (a wage--price dynamic). Even if the original shock fades, inflation can persist because expectations influence current contracts.

\section*{Part II: Income Inequality}

\subsection*{Question 5: Mean and median income}

Incomes: \(20{,}000,\ 30{,}000,\ 45{,}000,\ 70{,}000,\ 135{,}000\).

\subsubsection*{(1) Mean}

Sum:
\[
20{,}000+30{,}000=50{,}000
\]
\[
50{,}000+45{,}000=95{,}000
\]
\[
95{,}000+70{,}000=165{,}000
\]
\[
165{,}000+135{,}000=300{,}000
\]

Mean:
\[
\bar{y}=\frac{300{,}000}{5}=60{,}000
\]

\textbf{Answer:} Mean \(=\$60{,}000\).

\subsubsection*{(2) Median}

With 5 observations, the median is the 3rd income (ordered list):
\[
\text{Median}=45{,}000
\]

\textbf{Answer:} Median \(=\$45{,}000\).

\subsubsection*{(3) Why mean differs}

The high income (\$135,000) pulls the mean upward. The median is less sensitive to extremes, so it can better represent the “typical” household in a skewed distribution.

\newpage
\subsection*{Question 6: Income shares}

Total income \(=300{,}000\).

\subsubsection*{(1) Total income}

\[
Y=300{,}000
\]

\subsubsection*{(2) Share of top household}

Top household income \(=135{,}000\).

\[
\text{Top share}=\frac{135{,}000}{300{,}000}\times 100=45\%
\]

\textbf{Answer:} \(45\%\).

\subsubsection*{(3) Share of bottom two households}

Bottom two: \(20{,}000+30{,}000=50{,}000\).

\[
\text{Bottom two share}=\frac{50{,}000}{300{,}000}\times 100=16.67\%
\]

\textbf{Answer:} \(16.67\%\).

\subsubsection*{(4) Interpretation}

The top household receives nearly half of all income in this small economy, indicating high concentration. Bottom two receive only about one-sixth, highlighting inequality.

\subsection*{Question 7: Policy and redistribution}

Transfer \$8,000 to each of two lowest; tax \$16,000 on highest.

Original incomes:

A: 20,000

B: 30,000

C: 45,000

D: 70,000

E: 135,000
\newpage
\subsubsection*{(1) New incomes}

A:
\[
20{,}000+8{,}000=28{,}000
\]

B:
\[
30{,}000+8{,}000=38{,}000
\]

E:
\[
135{,}000-16{,}000=119{,}000
\]

C and D unchanged.

New set:

A: 28,000

B: 38,000

C: 45,000

D: 70,000

E: 119,000

\subsubsection*{(2) New mean}

Check total: transfers are funded by tax, so total income should remain 300,000. Sum to verify:
\[
28{,}000+38{,}000=66{,}000
\]
\[
66{,}000+45{,}000=111{,}000
\]
\[
111{,}000+70{,}000=181{,}000
\]
\[
181{,}000+119{,}000=300{,}000
\]

Mean:
\[
\bar{y}=\frac{300{,}000}{5}=60{,}000
\]

\textbf{Answer:} Mean remains \(\$60{,}000\).

\subsubsection*{(3) Inequality up or down}

Inequality decreases because the bottom rises and the top falls. For example, top-to-bottom ratio:

Before:
\[
135{,}000/20{,}000=6.75
\]

After:
\[
119{,}000/28{,}000\approx 4.25
\]

The distribution is less spread out.

\subsection*{Question 8: Poverty threshold \$28,000}

\subsubsection*{(1) Below poverty line}

Using original incomes:

20,000 $<$ 28,000 \(\rightarrow\) poor

30,000, 45,000, 70,000, 135,000 \(\rightarrow\) not poor

\textbf{Answer:} Only Household A is below.

\subsubsection*{(2) Poverty rate}

\[
\frac{1}{5}\times 100=20\%
\]

\textbf{Answer:} \(20\%\).

\subsubsection*{(3) One limitation}

A poverty line is a sharp cutoff: someone earning \$27,900 is “poor” while \$28,100 is “not poor,” though living standards are nearly identical. It also may not adjust well for local cost-of-living differences.

\section*{Part III: Housing Affordability and Real Estate Markets}

\subsection*{Question 9: Price-to-income ratio}

Income \(=95{,}000\), home price \(=760{,}000\).

\subsubsection*{(1) PTI}

\[
PTI=\frac{760{,}000}{95{,}000}=8
\]

\textbf{Answer:} \(8\).

\subsubsection*{(2) Interpretation}

A PTI of 8 means the home costs eight years of gross income -- typically considered severely unaffordable, especially once interest costs and down payments are considered.

\subsubsection*{(3) Why incomplete}

PTI ignores mortgage rates, down payment constraints, property taxes, condo fees, and household composition; also ignores whether the house is larger/better quality than typical.

\subsection*{Question 10: Mortgage payments and interest rates}

Loan \(P=500{,}000\), amortization 25 years \(=300\) months.

Monthly payment formula:
\[
M = P\cdot \frac{r(1+r)^n}{(1+r)^n-1}
\]
where \(r\) is monthly interest rate and \(n=300\).

We compute two cases.

\subsubsection*{Case 1: 2.5\% annual}

Monthly rate:
\[
r_1=\frac{0.025}{12}=0.00208333
\]

Compute \((1+r_1)^{300}\approx 1.867\) (approx). Then:
\[
M_1 \approx 500{,}000 \cdot \frac{0.00208333 \cdot 1.867}{1.867-1}
\]

Numerator: \(0.00208333\times 1.867 \approx 0.003889\)

Denominator: \(0.867\)

Fraction: \(0.003889/0.867 \approx 0.004484\)

\[
M_1 \approx 500{,}000 \times 0.004484 \approx 2{,}242
\]

So payment is about \$2,240/month.

\subsubsection*{Case 2: 5.5\% annual}

Monthly rate:
\[
r_2=\frac{0.055}{12}=0.00458333
\]

Compute \((1+r_2)^{300}\approx 3.94\) (approx).
\[
M_2 \approx 500{,}000 \cdot \frac{0.00458333 \cdot 3.94}{3.94-1}
\]

Numerator: \(0.00458333\times 3.94\approx 0.01806\)

Denominator: \(2.94\)

Fraction: \(0.01806/2.94\approx 0.00614\)

\[
M_2 \approx 500{,}000\times 0.00614 \approx 3{,}070
\]

So payment is about \$3,070/month.

\subsubsection*{(1) Change}

\[
\Delta M \approx 3{,}070-2{,}240 \approx 830
\]

\textbf{Answer:} Monthly payment rises by about \$800–\$850/month.

\subsubsection*{(2) Why affordability worsens even if prices stabilize}

Higher interest rates raise the monthly payment for any given loan size, and at renewal they raise cash-flow burdens. Even if the house price stops rising, households can still be “priced out” by monthly payments.

\subsection*{Question 11: Rental affordability}

Income \(=3{,}800\)/month. Rent rises from 1,200 to 1,850.

\subsubsection*{(1) Rent-to-income ratios}

Before:
\[
\frac{1{,}200}{3{,}800}=0.3158 \Rightarrow 31.58\%
\]

After:
\[
\frac{1{,}850}{3{,}800}=0.4868 \Rightarrow 48.68\%
\]

\textbf{Answer:} About 31.6\% \(\rightarrow\) 48.7\%.

\subsubsection*{(2) Threshold}

A common benchmark is 30\% (and sometimes 50\% for “severe burden”).

\subsubsection*{(3) Link to inequality}

Rising rents hit lower- and middle-income renters hardest because housing is a necessity. It reduces ability to save, invest in education, or accumulate wealth-widening the gap between renters and homeowners.

\subsection*{Question 12: Supply elasticity and prices (explanations)}

\subsubsection*{(1) Restrictive zoning}

Restrictive zoning makes supply less responsive (more inelastic). When demand rises, fewer new units are built, so prices and rents increase more.

\subsubsection*{(2) Increased immigration}

Raises population and household formation, increasing demand for rentals and ownership. With slow supply response, affordability worsens through higher rents/prices.

\subsubsection*{(3) New transit infrastructure}

Transit increases accessibility and raises desirability of nearby locations. It can increase housing demand in those areas (prices rise). Over time, if zoning allows densification around transit, supply can increase and moderate long-run price growth.

\subsubsection*{(4) Decline in mortgage interest rates}

Increases borrowing capacity; demand shifts outward; prices rise, especially when supply is constrained.

\section*{Part IV: Integrated Policy Analysis}

\subsection*{Question 13: Inflation and housing (combined mechanisms)}

High inflation often leads central banks to raise interest rates. Higher rates:

Slow house price growth by reducing mortgage borrowing capacity and cooling demand.

Increase mortgage payments for variable-rate borrowers immediately and for fixed-rate borrowers at renewal.

Raise rental demand because higher payments and stress tests price out marginal buyers, pushing them into renting, which can raise rents.

So, it is possible for house price inflation to slow while renter cost pressures intensify.

\subsection*{Question 14: Inequality and housing wealth}

Rising house prices increase wealth inequality because housing is a key asset.

Homeowners vs renters: Owners gain capital gains and build equity; renters do not. Higher prices also raise barriers to entry for renters to become owners.

Older vs younger households: Older cohorts are more likely to already own homes; they benefit from appreciation. Younger households face higher entry prices, larger down payments, and higher debt burdens.

This creates an intergenerational “asset divide.”



\end{document}